% Generated draft: requires mathematical and source review.
\documentclass[11pt]{article}
\usepackage[margin=25mm]{geometry}
\usepackage{fontspec}
\setmainfont{Latin Modern Roman}
\usepackage{amsmath,amssymb,mathtools,unicode-math}
\setmathfont{Latin Modern Math}
\usepackage{enumitem,needspace,xurl,hyperref}
\setlength{\parindent}{0pt}
\setlength{\parskip}{5pt}
\setlength{\emergencystretch}{3em}
\newcommand{\N}{\mathbb{N}}
\newcommand{\Z}{\mathbb{Z}}
\newcommand{\Q}{\mathbb{Q}}
\newcommand{\R}{\mathbb{R}}
\newcommand{\C}{\mathbb{C}}
\newcommand{\F}{\mathbb{F}}
\newcommand{\A}{\mathbb{A}}
\newcommand{\PP}{\mathbb{P}}
\newcommand{\E}{\mathbb{E}}
\newcommand{\eps}{\varepsilon}
\newcommand{\dd}{\,\mathrm{d}}
\begin{document}
\section*{UnsolvedMath ID: 20001655}
\textbf{Source problem code:} catalogue entry
\subsection{Definitions and mathematical statement}
For a marked minimal \(F_N\)-tree \(T\), write \(\|g\|_T=\inf_xd(x,gx)\). Its dual algebraic lamination, in the convention relevant here, is
\[
 L^2(T)=\bigcap_{\epsilon>0}\overline{\{(g^{-\infty},g^{+\infty}):1\ne g\in F_N,\ \|g\|_T<\epsilon\}}
 \ \subseteq\partial^2F_N,
\]
where \(\partial^2F_N=(\partial F_N)^2\setminus\operatorname{diag}\), closures are in that space, and opposite orientations are included. The source asks to understand the ``GL-theory of fine decomposition'' through dual laminations in both directions and through systems of partial isometries. It supplies no reference or precise target equivalence for ``GL-theory''.

The attached aid treats a restricted graph-of-actions class. Assume
\[
 F_N=H_1*\cdots*H_k*F_r
\]
comes from a finite-quotient graph of actions with trivial edge groups and strictly positive simplicial edge lengths, and all nondegenerate vertex trees \(Y_i\) are free, minimal, indecomposable dense-orbit \(H_i\)-trees. Free means trivial point stabilizers. Indecomposable means that any nondegenerate arc is covered by finitely many translates of any other nondegenerate arc with consecutive nondegenerate overlaps. Each \(\partial^2H_i\) is viewed inside \(\partial^2F_N\).

For these assumptions the aid reports
\[
 L^2(T)=\bigsqcup_{i=1}^k\ \bigcup_{gH_i\in F_N/H_i}gL^2(Y_i).
\]
Distinct induced component laminations are disjoint. It reports that the minimal nonempty closed invariant sublamination blocks have free-factor supports exactly the conjugacy classes \([H_i]\). Free-factor support means the smallest free-factor conjugacy class carrying the block's leaves. Consequently the lamination recovers the nondegenerate vertex-group orbits and each component tree up to its observers topology, the topology on its metric completion together with its ends generated by tree directions. It does not recover component metrics, simplicial edge lengths or attachment points. The result is restricted to the stated free indecomposable dense-orbit class, not arbitrary fine decompositions.

The aid proposes compatible length measures and compact-heart systems of partial isometries as extra data for metric reconstruction. A system of partial isometries records isometries between specified subtrees of a compact tree; its domains, range identifications, length data and the skeleton's edge lengths and attachment choices must be specified. The supplied summary does not furnish a complete algorithm or a general uniqueness theorem for that enhancement.

The source also records Cormac Walsh's discussion of the horofunction boundary of outer space \(X_N\) for the asymmetric Lipschitz metric. In the usual volume-one marked-graph model, \(d(x,y)\) is the logarithm of the least Lipschitz constant of a marking-compatible map. The standard horofunction construction uses, for a basepoint \(x_0\),
\[
 h_x(z)=d(z,x)-d(x_0,x),
\]
with uniform convergence on compact subsets, equivalently continuous functions modulo additive constants. The source instead prints ``homotopy to a constant'' and uses \(C^0\)-like notation. Homotopy is not that quotient: real-valued continuous functions are all homotopic to constants, and distance functions need not vanish at infinity. Thus the displayed standard interpretation is a labeled repair, not a verbatim valid definition. The source's statement that this boundary is coarser than the usual boundary is retained as contextual attribution, requiring a specified comparison of compactifications and metric direction. The supplied component theorem does not resolve that boundary comparison.
\subsection{Short English statement}
Relate fine graph-of-actions decompositions to dual laminations and partial isometries. In the supplied restricted class, lamination blocks recover vertex-group supports and observers-topological components, while metrics and gluing require extra data. The separate horofunction paragraph needs additive constants in place of its printed homotopy quotient.
\subsection{Sources}
\begin{itemize}
\item AIM workshop \emph{The geometry of the outer automorphism group of a free group}, Problem 5.2 and the accompanying horofunction discussion attributed to Cormac Walsh: \url{https://aimath.org/WWN/outerauto/outerauto.pdf}.
\item Local registry \texttt{data/unsolvedmath/problems.json}, permanent ID 20001655, source code \texttt{catalogue entry}; canonical source \texttt{aim-geometric-group-theory-notes.json}, \texttt{notes[163]}.
\item The source supplies only the abbreviation ``GL-theory'' and the Walsh attribution, with no full reference for fine decomposition, lamination reconstruction or the compactification comparison.
\item Attached machine-generated attempt 1: partial result, run status partially solved, internal rigor label ``proved'', low novelty confidence; not peer reviewed and requiring independent verification before citation.
\end{itemize}
\end{document}
